% ======================================================================
% 40_discussion_conclusion.tex — long Discussion + Conclusion fragment.
% Table floats T1-T5 live in tables_main.tex (tab:FlxPinGof,
% tab:ExtPinGof, tab:AdoptedXi, tab:Redesign, tab:BiPulley) and are
% referenced here, not duplicated. The one new float is
% fig:ExtensorRoute_redesigned (file to be staged as
% Figures/ExtensorRoute_redesigned.pdf per plan section 6).
% Uses \diameter (requires mathabx in the preamble, flagged by the
% plan's cleanup pass). Numbers of record per results_ledger.md and
% chapters/30-torque.tex / 94-AppendixC.tex (committed 0d23da03).
% ======================================================================

\section{Discussion}\label{sec:discussion}

\subsection{Sources of Error in Isometric Torque Generation}\label{sec:disc_error}

The improved BPA force characterization removes part of the torque-prediction error, but only part (Figure~\ref{fig:FlxPin_group}\textbf{A,B}). The remainder traces to the test system itself. The isometric system is not rigid: the winch cable, the artificial tendons, and the attachment brackets all bend and stretch under load, and the correction family $\chi_{0}$, $\chi_{1}$, and $\chi_{2}$ was created to represent them. Positioning tolerances, imperfect path modeling, and variation between individual actuators contribute as well.

The modeling of extensor paths that wrap around the knee is an approximation with inherent error, and the test-stand geometry amplifies it. When the moment arm is short, small changes to it have a large effect on the resultant torque \citep{young_analyzing_2019}, which is why moment-arm determination receives so much attention in musculoskeletal modeling \citep{brand_model_1982, menegaldo_moment_2004, suderman_moving_2012}. We chose to wrap the BPA around the joint rather than add a patella with a deterministically designed location and friction; the resulting path, which passes between the femoral condyles close to the migrating center of rotation of the biomimetic knee \citep{steele_biomimetic_2018}, resists any straight-line idealization (Supplementary Figure~\ref{fig:passThru}). At high flexion, the shorter actuators were also observed to stretch and compress (Figure~\ref{fig:testJigs}\textbf{C}), and the same loading occasionally slipped the actuator off the anterior bolt head that holds the femur to the knee, displacing the path by up to $\pm$\qty{20}{\mm} in the $Z$ direction and changing both the force vector and the torque.

The pinned extensor tests exposed one further error source, with a signature distinct from compliance. Figure~\ref{fig:ExtPin_group} shows a deviation between measured and originally predicted torque that widens as the expected torque grows, yet the hybrid calculation matches or exceeds the prediction except at the highest flexion angles for the $l_{0}=\qty{41.5}{\cm}$ actuator. That combination points to a loss of functional muscle length rather than excess compliance, and it is why the wrapping-loss term $\chi_{3}$ was created in Equation~(\ref{eq:deltaL}). Bent over the joint, the actuator behaves like a beam, and the wrapped arc stops contributing usable contraction, so the actuator behaves as if it were more contracted than it physically is; this is the fictive strain of Section~\ref{sec:fictive_contraction}, and Section~\ref{sec:adopted_identification} quantifies it.

Working from the simple pinned knee to the complex biomimetic one isolated each of these mechanisms in turn. The products of that progression, a predictor for bending and axial compliance, a constant length offset, and a measure of the usable length lost as the actuator bends around a joint, are what make first-iteration torque design tractable for BPA-driven joints.

\subsection{Two-Stage Re-Identification and the Adopted Correction Terms}\label{sec:adopted_identification}

The coefficients quoted in the Results belong to the initial single-bracket identification (Section~\ref{sec:identification}). Re-identifying the correction terms over the full test sets, with the two-bracket flexor model and the locked extensor refit, yields the adopted values: $\chi_{0} = \qty{+3.94}{\mm}$, $\chi_{1} = \qty{3.998e4}{\N\per\m}$, and $\chi_{2} = \qty{1.473e4}{\N\per\m}$ for the flexor, and $\chi_{0} = \qty{-6.38}{\mm}$ with $\chi_{3} = 0.158$ for the extensor on the same locked stiffness pair (Table~\ref{tab:AdoptedXi}).

Each term has a direct physical reading. The offsets $\chi_{0}$ correct the effective muscle length by a constant, and the flexor and extensor offsets carry opposite signs, as the initial identification already showed. The stiffness pair is an effective system stiffness rather than the beam stiffness of any single part: it lumps the two attachment brackets together with the fixtures and the winch cable, which is why one pair can serve both configurations once the extensor refit locks it. The wrapping factor $\chi_{3}$ scales the usable length that fictive strain removes, $\Delta l = \chi_{3}\, R\, \abs{\theta_{\text{wrap}} - \theta_{\text{k}}}\, \zeta^{2}$ with $\zeta = 1 - \varepsilon^{*}$ (Equations~(\ref{eq:deltaL}) and (\ref{eq:comp})). Because $\varepsilon^{*}$ is normalized by the unloaded \qty{620}{\kPa} contraction $\varepsilon_{620}$ of the free actuator, the loss is expressed in fictive-contraction units: it vanishes as the actuator approaches full contraction, where $\zeta \to 0$, and grows quadratically toward resting length, exactly where a wrapped actuator has the least contraction to spare.

The pinned-flexor data do not uniquely determine the stiffness pair. A family of pairs reproduces the training tests nearly equally well, so the adopted pair was selected by a requirement external to the fit: both route-redesign configurations must meet the required \qty{5}{\percent} torque-overshoot target simultaneously. Under that rule, the flexor redesign closed at a minimum margin of $+5.00$~\% and the locked extensor redesign at $+5.043$~\% at a knee angle of \ang{-20.2}; the finalized redesigns hold $+5.01$~\% and $+5.02$~\% (Section~\ref{sec:route_redesign_results}).

Tables~\ref{tab:FlxPinGof} and \ref{tab:ExtPinGof} quantify what the corrections buy. On the pinned-flexor set, the corrected model reduces the training-mean RMSE from \qty{8.12}{\newton\meter} to \qty{1.56}{\newton\meter} under the zero-correction prediction, the same kinematics and force model with every correction term disabled, where the mean averages the four training tests; it holds the FVU at or below 0.17 on all five tests and validates on the held-out $l_{0}=\qty{47.9}{\cm}$ test at \qty{2.29}{\newton\meter}. On the extensor pool, the corrected model improves the RMSE of every test, roughly halving the fitting-set mean from \qty{2.60}{\newton\meter} to \qty{1.16}{\newton\meter}, with the held-out $l_{0}=\qty{45.7}{\cm}$ test at \qty{0.49}{\newton\meter}; the FVU stays below unity on every test except the $l_{0}=\qty{43.2}{\cm}$ tendon test. The biomimetic extensor test carries two goodness-of-fit pairs in this paper because it was evaluated at two parameter sets: the Results quote the initial extensor terms, and at the adopted terms the corrected prediction improves the RMSE from \qty{3.03}{\newton\meter} to \qty{1.17}{\newton\meter} and the FVU from 4.49 to 0.67 against the rigid-body baseline (Figure~\ref{fig:Ext10mm_52cm}), so the shared-pair model explains the bulk of that configuration's measured torque variance as well as its level.

The solver assignment follows the problem. The identification retains the genetic multiobjective solver, an elitist nondominated-sorting genetic algorithm, because the Pareto front it returns is what the holdout-validation procedure selects from. The route redesign is served by the surrogate-plus-pattern-search pipeline of Section~\ref{sec:route_redesign}, a surrogate-assisted global search refined by a derivative-free mesh-adaptive pattern search, because a single scalar cost is optimized there and each route evaluation is expensive.

The corrected force balance also admits an effective spring--damper reading in relative strain. Writing the force model as $F(\varepsilon^{*}, P^{*})$, its local stiffness is $k(\varepsilon^{*}, P^{*}) = \left.\partial F/\partial\varepsilon^{*}\right|_{P^{*}}$, and setting the mass and the external force of that analogy to zero leaves an effective damping $c = k\,\varepsilon^{*}/\dot{\varepsilon}^{*}$, the ratio between the restoring force and the strain rate that the identification balances implicitly. This reading connects the quasistatic corrections to the dynamic regime the robot leg will operate in, where the same compliance sets the joint's passive damping.

\subsection{Route Redesign Meeting Human Isometric Torque Targets}\label{sec:route_redesign_results}

The redesign addresses the single-actuator shortfall documented in the Results: one BPA per route does not reach the human flexor envelope through the mid-flexion range. Every redesigned configuration therefore carries two \diameter\qty{20}{\mm} BPAs per route, a \qty{620}{\kPa} pressure ceiling, and a maximum usable contraction of 0.255 in the force model, and every family was produced by the procedure of Section~\ref{sec:route_redesign} under the adopted correction terms of Table~\ref{tab:AdoptedXi}. The targets are maximum-isometric torque profiles computed on the Gait2392 lower-extremity model \citep{delp_opensim_2007, seth_opensim_2018}, the vastus medialis for the extensor and the biceps femoris short head for the flexor, and the redesign itself runs on the robot-body model derived from it. each enlarged by a required \qty{5}{\percent} overshoot margin so that a route merely touching the target at any angle is rejected. Knee angles are negative in flexion throughout this section, and Table~\ref{tab:Redesign} summarizes the four families.

The direct extensor route exceeds the human envelope in magnitude at every angle, with zero shortfall points. Across the design range, the predicted two-BPA torque spans \qty{35.45}{\newton\meter} to \qty{109.4}{\newton\meter} against a human-scaled requirement of \qty{20.25}{\newton\meter} to \qty{50.32}{\newton\meter}, and the required overshoot is consumed as the thinnest margin exactly as designed, $+5.02$~\% at the worst-case angle of \ang{-20.2}. The route respects the wrap-release schedule of the route-elimination rule, releasing its six sequential contacts at knee angles of \ang{-89.8}, \ang{-55.7}, \ang{-47.8}, \ang{-20.2}, \ang{6.1}, and \ang{7.4}, and the optimized design places the femur-side origin at $(40.0, 137.0, 0)$~mm and the tibial insertion at $(34.4, -197.1, 0)$~mm, with a resting length of \qty{827.7}{\mm} and \qty{36.8}{\mm} of tendon (Figure~\ref{fig:ExtensorRoute_redesigned}).

\begin{figure}[htbp]
\begin{center}
\includegraphics[width=1.0\textwidth]{Figures/ExtensorRoute_redesigned.pdf}
\end{center}
\caption{Route geometry of the redesigned extensor route, drawn in the femur frame. Eight pose panels sweep the design range from full flexion ($\theta_\text{k} = \ang{-120}$) through the frame immediately before each of the six contact-elimination events to full extension ($\theta_\text{k} = \ang{10}$); knee angle is negative in flexion. Each panel draws the active route over the femur-cylinder, femur-wall, condyle, and tibia clearance envelopes of the route-elimination rule, with the local bend-radius checks dashed. The square marker is the femur-side origin and the diamond marker is the tibial insertion.}\label{fig:ExtensorRoute_redesigned}
\end{figure}

The extensor route was then redesigned once more by allowing the tibial-side route points to trace a cam profile, and this cam redesign is the current extensor design. The cam converts the razor-thin margin of the direct route into genuine headroom. On the final geometry, which folds the optimizer winner together with a verified adjustment of the distal route point and tendon length, the minimum torque margin is $+30.83$~\% with respect to the required envelope and $+39.77$~\% with respect to the plain human target; the unadjusted optimizer solution carried $+31.30$~\% and $+42.34$~\% at the same checks. The pressure required to meet the envelope spans \qty{269}{\kPa} to \qty{507}{\kPa} over the range of motion, comfortably above the \qty{150}{\kPa} regulated floor of the valve manifold that supplies the hardware, so the remaining pressure is a design resource. The geometry confirms that the headroom is not bought with routing risk: the tibial wrap follows a bend radius of \qty{96.9}{\mm}, the guiding route point clears the femoral-condyle envelope by \qty{43.1}{\mm} at \ang{10} of knee extension, and the moment arm spans \qty{59.4}{\mm} to \qty{102.8}{\mm} with a mean of \qty{72.4}{\mm}.

The flexor route, built on the same two-BPA arrangement sharing one attachment route, meets the flexor target with a minimum margin fraction of 0.0501 against the required 0.05. Its predicted knee-axis torque spans \qty{-81.23}{\newton\meter} to \qty{-3.864}{\newton\meter} against the human-scaled requirement of \qty{-23.16}{\newton\meter} to \qty{-2.461}{\newton\meter} (peak magnitude \qty{23.16}{\newton\meter}), again exceeding the envelope at every angle. The flexor prediction applies the adopted correction terms, including the wrapping factor identified on the extensor configuration.

Allowing the artificial tendon to be reeved over pulleys before insertion extends the search space, and the knee-flexor run quantifies what that freedom is worth. The mono-pulley variant converged to a collision-feasible layout with zero bound violations, reeving both actuators over a single tackle-line part with the pulley gain fixed at 1:1, so the pulley acts as a redirect rather than a force multiplier. Even so, the redirect alone deepened the peak predicted flexion torque from \qty{81.23}{\newton\meter} on the direct route to \qty{85.74}{\newton\meter} on the same knee-axis measure, \qty{85.56}{\newton\meter} on the insertion-side measure after the tackle's force loss, a gain of \qty{4.51}{\newton\meter} or \qty{5.6}{\percent}, while holding the off-axis torque to at most \qty{0.803}{\newton\meter}, under \qty{1}{\percent} of the peak. The direct and pulley runs optimize different objective functions, so their objective values, 74.47 and 0.08465, are not commensurable; the torque comparison, not the objective comparison, is the meaningful one. The predicted peaks of the direct and pulley routes exceed the human flexor scale of \qty{23.16}{\newton\meter} by factors of 3.51 and 3.70.

The bi-pulley sweep establishes where that freedom runs out on anatomical routes. Applied to five muscle sites of the robot-body model's leg, all five runs failed the torque requirement (Table~\ref{tab:BiPulley}): every run closed with a nonpositive torque margin, and the medial-gastrocnemius and biceps-femoris-short-head runs terminated at the infeasibility sentinel. For the four runs that report a finite route-to-route clearance, the binding constraint was geometry rather than torque: the minimum pair clearances spanned \qty{4.26}{\mm} to \qty{19.14}{\mm} against the \qty{25}{\mm} separation required between adjacent routes, the sum of two \diameter\qty{10}{\mm} actuator radii and a \qty{5}{\mm} bundle margin. For the biceps femoris short head site, whose target is \qty{40.2}{\newton\meter} at a maximum isometric force of \qty{804}{\N}, the search exited with an infeasibility flag. The practical outcome is a screening rule: reeved routes do not yet fit these corridors, and design effort is better spent on clearance-relaxed attachment sites.

Physical validation of the redesigned extensor route has begun. In the first test of the \diameter\qty{20}{\mm} extensor on the rebuilt route, three isometric points were held near full extension: measured knee torque of \qty{1.163}{\newton\meter} at \ang{10}, \qty{7.849}{\newton\meter} at \ang{11}, and \qty{23.437}{\newton\meter} at \ang{10}, against hung loads of \qtylist{3.7;25.2;76}{\newton} at gauge pressures of \qtylist{388.5;388;592}{\kPa}. The first and third points sit at essentially the same knee angle and span a twentyfold range of held torque, so they begin to map the route's response near extension as pressure and load vary. The full angle sweep, and the corresponding flexor measurements, are pending.

The same pipeline was extended from the knee to all 27 muscle-actuators of the robot-body model, and we report that extension as a status rather than a completed result. The first pass reached 19 of the 27 actuators, and its chosen routes showed the penalty-based formulation to be the binding limitation: all but one violated the feasibility constraints, the exception being the biceps femoris long-head pair at a worst-angle margin of $+0.03$~\%, indistinguishable from zero. The formulation was then restructured so that the task envelope, the joint-angle window each route must serve, acts as a hard feasible set rather than a penalty. In the restructured pass, twelve actuators terminated on a concatenation error in the route builder, leaving fifteen finalized routes; eight were chosen from constraint-feasible configurations (semimembranosus, biceps femoris long head, tensor fasciae latae, soleus, lateral gastrocnemius, tibialis anterior, iliacus, and psoas), and six of those met the torque requirement at every angle, with worst-angle margins from $+0.25$~\% (biceps femoris long head) to $+250.7$~\% (tensor fasciae latae). All of these redesign solutions remain optimizer-stage results: the reference model contains no pulley entities, no via points or markers were added to it for these routes, and the redesigned geometry has not yet been propagated into the dynamic simulation models.

\subsection{Design Factors for Isometric Torque}\label{sec:design_factors}

Read as a design problem, isometric torque about a BPA-actuated joint factors into three quantities the designer controls: the moment arm $r$, the normalized contraction $\varepsilon^{*}$ at which the actuator operates, and the compliance of the path and its fixtures, carried by the $\chi$ terms. Each result in this paper maps onto one corner of that triangle. Short moment arms amplify every error, so stiffness must be budgeted where the arm is short \citep{young_analyzing_2019}. The route should keep $\varepsilon^{*}$ high through the functional range, because the wrapping loss of Equation~(\ref{eq:deltaL}) scales with $(1 - \varepsilon^{*})^{2}$ and therefore grows exactly where an under-contracted actuator is weakest. Moment arm and contraction budget trade against each other, and the cam redesign shows the trade can be favorable: it buys \qty{59.4}{\mm} to \qty{102.8}{\mm} of moment arm while improving both margin and clearance. Pressure headroom, finally, is a design resource: the cam route meets its envelope between \qty{269}{\kPa} and \qty{507}{\kPa} against a \qty{620}{\kPa} ceiling.

These factors situate our work among musculoskeletal robot legs. Multifilament muscle routing has produced human-scale lower limbs \citep{kurumaya_musculoskeletal_2016}, anthropomorphic arms \citep{hitzmann_anthropomorphic_2018}, and full human-mimetic humanoids \citep{asano_musculoskeletal_2019}, and tendon-driven bodies have supported investigations from reflexive hopping \citep{liu_robotic_2018} to human-like manipulation \citep{chen_realizing_2019} and cognitive locomotion \citep{schilling_adaptive_2022}. Linkage-based knees manage the moment-arm problem geometrically \citep{romero_four-bar_2020}, reviews of dynamic multilegged robots treat actuation routing as a first-order design decision \citep{he_mechanism_2020}, and tunable-compliance actuators make fixture compliance a design element rather than an error source \citep{leibach_development_2020}. What these designs share with ours is the recognition that the attachment route, not the actuator alone, sets the deliverable torque. What the corrected model adds is a quantitative accounting of the compliance and wrapping losses along that route, so that a route can be designed against a torque target rather than checked after the fact. \citet{morrow_optimization_2020} optimized artificial-muscle placement on a bipedal robot with a rigid-body torque model, and the correction terms quantified here are the physics such rigid-body placements leave unmodeled.

The correction family also complements the BPA force-modeling literature. Static models of the actuator in isolation, from the classical McKibben analyses \citep{chou_measurement_1996} through approximation formulas \citep{sarosi_new_2012} to geometrically grounded models of the Festo muscle \citep{martens_modeling_2017, tondu_static_2022}, describe the actuator itself; none of them carries attachment compliance, length offsets, or wrapping losses into a joint-torque prediction. The terms validated here attach to the force characterization of \citet{bolen_2026} and travel with it across joint configurations under one consistent parameter set.

\subsection{Limitations and Future Work}\label{sec:limitations}

The redesign results carry caveats that bound their use. The flexor redesign assumes both actuators share one attachment route, so their combined force acts on a single moment arm; making the mirrored pair explicit is the next design refinement, and it will relax that assumption. All redesign solutions are optimizer-stage: they have not been propagated into the simulation models, and the physical validation rests on three torque points near full extension, with the full sweep and the flexor measurements pending. The stiffness pair itself was selected by an external design requirement rather than uniquely identified, because the pinned-flexor data admit a family of near-equivalent pairs; an independent stiffness measurement of the test stands would close that non-identifiability and either confirm or replace the adopted pair. The wrapping factor $\chi_{3}$ was identified on the \diameter\qty{10}{\mm} extensor family and applied unchanged to the \diameter\qty{20}{\mm} redesigns, so its transfer across actuator diameters is untested. On the characterization side, the fictive-strain formulation inherits the conclusion of our force characterization that the unloaded \qty{620}{\kPa} contraction $\varepsilon_{620}$ should be recorded for each actuator used on a robot \citep{bolen_2026}, since every correction term is normalized by it. The whole-body extension identified its own next steps: the route-builder concatenation failure needs a fix before the twelve crashed routes can be rebuilt, and the two torque-short feasible routes need either larger actuators or relaxed corridors. One scope statement, finally, belongs with these caveats: the correction terms are identified for these test stands and this actuator family, and other joints or actuator sizes should repeat the two-stage procedure rather than reuse the values. That procedure is inexpensive next to the test campaign it saves, which is the point of having it.

\section{Conclusion}\label{sec:conclusion}

This work asked what a braided pneumatic actuator actually delivers about a joint, and answered with a corrected torque model and the design workflow built on it. Isometric knee-torque measurements on pinned and biomimetic test stands, interpreted through a hybrid calculation that substitutes the measured muscle length and moment arm, isolated the error sources that rigid-body torque models miss: attachment compliance, a constant effective-length offset, and a wrapping loss that acts like fictive strain, expressed through the unloaded \qty{620}{\kPa} contraction of the free actuator. A two-stage identification recovered the correction terms with one stiffness pair shared across the flexor and extensor configurations, roughly halving fitting-set error, cutting held-out error by two thirds or more, and carrying the corrected model from the pinned knee to the biomimetic knee without refitting. The attachment routes themselves were then redesigned under the corrected model: two \diameter\qty{20}{\mm} actuators per route, graded against OpenSim human isometric torque profiles with a required \qty{5}{\percent} overshoot, meet or exceed the human flexor and extensor envelopes at every angle, the cam extensor design converts the thinnest margins into design headroom, the pulley sweep screens where reeved routing runs out of corridor, and the first physical torque points of the redesigned extensor route have been measured near full extension. The design-factor reading is the practical summary: moment arm, normalized contraction, and path compliance are the three quantities a designer of BPA-actuated joints controls, and each now carries a measurable, correctable model. First-iteration torque design for such joints is, as a result, no longer a guess.
